\documentclass{article}
\usepackage{amsmath,amsthm,amssymb,amsfonts}
\usepackage{fullpage}
\usepackage{times}
\usepackage{verbatim}
\usepackage{graphicx}
\usepackage{float}
\usepackage{comment}
\usepackage{color}
\usepackage{minted}
\usepackage{appendix}
\usepackage{subcaption}
\usepackage{svg}
\usepackage{pdfpages}
\usepackage{hyperref}
\usepackage{mathtools}
\usepackage{gensymb}
% \usepackage{minipage}
\usepackage[varqu]{zi4} % "Inconsolata" font - very readable for code
\usepackage[T1]{fontenc} % Ensures the font renders correctly

\usepackage{algorithm}
\usepackage{algorithmic}


\usepackage[most]{tcolorbox}

% Define the "theorem" environment 
\newtcbtheorem[auto counter, number within=section]{theorem}{Theorem}%
{
  enhanced,
  colback=blue!5!white,      % Very light blue background
  colframe=blue!75!black,    % Dark blue frame
  fonttitle=\bfseries,
  coltitle=white,
  boxrule=0.5mm,
  attach boxed title to top left={xshift=5mm, yshift*=-\tcboxedtitleheight/2},
  boxed title style={colback=blue!75!black}
}{thm} % "thm" is the prefix for referencing (e.g., \ref{thm:my-theorem})

% Define the "lemma" environment
\newtcbtheorem[auto counter, number within=section]{lemma}{Lemma}%
{
  enhanced,                % Enable advanced features
  colback=blue!5!white,    % Background color (very light blue)
  colframe=blue!75!black,  % Frame color (dark blue)
  fonttitle=\bfseries,     % Title font
  coltitle=white,          % Title text color
  boxrule=0.5mm,           % Thickness of the box border
  attach boxed title to top left={xshift=5mm, yshift*=-\tcboxedtitleheight/2},
  boxed title style={colback=blue!75!black}
}{lem} % "lem" is the prefix for referencing (e.g., \ref{lem:my-lemma})

% Define the "definition" environment
\newtcbtheorem[auto counter, number within=section]{definition}{Definition}%
{
  enhanced,
  colback=teal!5!white,      % Very light teal background
  colframe=teal!75!black,    % Dark teal frame
  fonttitle=\bfseries,
  coltitle=white,
  boxrule=0.5mm,
  attach boxed title to top left={xshift=5mm, yshift*=-\tcboxedtitleheight/2},
  boxed title style={colback=teal!75!black}
}{def} % "def" is the prefix for referencing (e.g., \ref{def}{my-definition})

% Define the "note" environment
\newtcolorbox{note}{
    colback=yellow!10!white,
    colframe=yellow!50!black,
    title=Note,
    fonttitle=\bfseries
}

% Define the "axiom" environment
\newtcbtheorem[auto counter, number within=section]{axiom}{Axiom}%
{
  enhanced,
  colback=violet!5!white,    % Very light purple background
  colframe=violet!75!black,  % Dark purple frame
  fonttitle=\bfseries,
  coltitle=white,
  boxrule=0.5mm,
  attach boxed title to top left={xshift=5mm, yshift*=-\tcboxedtitleheight/2},
  boxed title style={colback=violet!75!black}
}{ax}

% Define a proof box
\newtcolorbox{boxproof}{
  enhanced,
  breakable,                % Allows proof to break across pages
  colback=white,            % White background
  colframe=gray!50,         % Light gray frame (subtle)
  title={\textit{Proof.}},  % Title looks like standard proof
  coltitle=black,           % Black title text
  fonttitle=\bfseries,
  attach boxed title to top left={xshift=5mm, yshift*=-\tcboxedtitleheight/2},
  boxed title style={colback=white, frame hidden}, % Title floats in whitespace
  after upper=\hfill\qedsymbol % Puts the QED square at the end automatically
}

% Define a simple black box environment (No Title)
\newtcolorbox{answer}{
  enhanced,
  breakable,                % Allows box to break across pages
  colback=white,            % White background
  colframe=black,           % Black border
  boxrule=0.5mm,            % Matches the thickness of your other boxes
  coltext=black             % Ensures text inside is black
}

% Load the minted library for tcolorbox
\tcbuselibrary{minted} 

% --- 1. C++ CODE BOX (2 spaces per tab, No Line Numbers) ---
\newtcblisting[auto counter, number within=section]{cppcode}[2][]{%
  enhanced,
  colback=gray!5!white,        % Light gray background
  colframe=red!75!black,       % Red frame
  listing only,
  minted language=cpp,
  minted options={
      linenos=false,           % NO Line numbers
      fontsize=\small,         
      tabsize=2,               % 1 Tab = 2 Spaces
      breaklines=true,
      autogobble=true,
      fontfamily=tt            % Uses the new Inconsolata font
  },
  fonttitle=\bfseries,
  coltitle=white,
  title={Listing \thetcbcounter: #2}, 
  attach boxed title to top left={xshift=5mm, yshift*=-\tcboxedtitleheight/2},
  boxed title style={colback=red!75!black},
  #1
}

% --- 2. PYTHON CODE BOX (4 spaces per tab, No Line Numbers) ---
\newtcblisting[auto counter, number within=section]{pycode}[2][]{%
  enhanced,
  colback=gray!5!white,
  colframe=red!75!black,
  listing only,
  minted language=python,
  minted options={
      linenos=false,           % NO Line numbers
      fontsize=\small,
      tabsize=4,               % 1 Tab = 4 Spaces
      breaklines=true,
      autogobble=true,
      fontfamily=tt            % Uses the new Inconsolata font
  },
  fonttitle=\bfseries,
  coltitle=white,
  title={Listing \thetcbcounter: #2},
  attach boxed title to top left={xshift=5mm, yshift*=-\tcboxedtitleheight/2},
  boxed title style={colback=red!75!black},
  #1
}

\newcommand{\abs}[1]{\left| #1 \right|}
\newcommand{\norm}[1]{\lVert  #1 \rVert}
\newcommand{\parens}[1]{\left( #1 \right)}
\newcommand{\braces}[1]{\left\{ #1 \right\}}
\newcommand{\brackets}[1]{\left[ #1 \right]}
\newcommand{\angles}[1]{\langle #1 \rangle}

\setlength{\parindent}{0pt}

\newenvironment{code}{\captionsetup{type=listing}}{}

\title{Default Template}
\author{Eappen Nelluvelil}
\date{October 2026}

\begin{document}

\maketitle

\section{Analytical algorithm}

The universal DSA scheme is given below.
Assume $x \in \brackets{0, x_{0}}$ and that there are $G$ groups.
The $S_{N}$ equations are below:
\begin{align}
  \mu \frac{\partial \psi_{1}^{\ell+\frac{1}{2}}}{\partial x} + \sigma_{t, 1} \psi_{1}^{\ell+\frac{1}{2}} & = \frac{1}{4\pi} \sigma_{1 \rightarrow 1} \phi_{1}^{\ell},                                                                                                                \\
  \mu \frac{\partial \psi_{g}^{\ell+\frac{1}{2}}}{\partial x} + \sigma_{t, g} \psi_{g}^{\ell+\frac{1}{2}} & = \frac{1}{4 \pi} \sigma_{g \rightarrow g} \phi_{g}^{\ell} + \frac{1}{4 \pi} \sigma_{\parens{g - 1} \rightarrow g} \phi_{\parens{g-1}}^{\ell}, \quad g = 2, 3, \ldots, G.
\end{align}

The boundary conditions are below:
\begin{align}
  \psi_{g} \parens{0, \mu}     & = 0, \quad \mu > 0, \quad g = 1, 2, \ldots, G, \\
  \psi_{1} \parens{x_{0}, \mu} & = 1, \quad \mu < 0,                            \\
  \psi_{g} \parens{x_{0}, \mu} & = 0, \quad \mu < 0, \quad g = 2, 3, \ldots, G.
\end{align}

The diffusion acceleration step is given below:
\begin{align}
  -\frac{\partial}{\partial x} \frac{1}{3 \sigma_{t, 1}} \frac{\partial \phi_{1}^{\ell+1}}{\partial x} + \sigma_{R, 1} \phi_{1}^{\ell+1} & = \frac{\partial}{\partial x} \parens{J_{1}^{\ell + \frac{1}{2}} + \frac{1}{3 \sigma_{t, 1}} \frac{\partial \phi_{1}^{\ell + \frac{1}{2}}}{\partial x}},                                                                                                 \\
  -\frac{\partial}{\partial x} \frac{1}{3 \sigma_{t, g}} \frac{\partial \phi_{g}^{\ell+1}}{\partial x} + \sigma_{R, g} \phi_{g}^{\ell+1} & = \sigma_{\parens{g - 1} \rightarrow g} \phi_{\parens{g-1}}^{\ell+1} + \frac{\partial}{\partial x} \parens{J_{g}^{\ell + \frac{1}{2}} + \frac{1}{3 \sigma_{t, g}} \frac{\partial \phi_{g}^{\ell + \frac{1}{2}}}{\partial x}}, \quad g = 2, 3, \ldots, G.
\end{align}

Assume vacuum boundary conditions.
The cross-sections are defined as follows:
\begin{align}
  \sigma_{t, g} & = \sigma_{a, g} + \sigma_{g \rightarrow g} + \sigma_{g \rightarrow \parens{g + 1}}, \quad g = 1, 2, \ldots, G - 1, \\
  \sigma_{R, g} & = \sigma_{a, g} + \sigma_{g \rightarrow \parens{g + 1}}, \quad g = 1, 2, \ldots, G - 1,                            \\
  \sigma_{t, G} & = \sigma_{a, G} + \sigma_{G \rightarrow G},                                                                        \\
  \sigma_{R, G} & = \sigma_{a, G}.
\end{align}

We can just use diamond-differencing for both the $S_{N}$ and diffusion equations (look back at the NUEN 625 notes).
We can also define a routine to define the cross-sections.

For the time being, take $x_{0} = 10$ cm and the following as the cross-sections for all groups $g = 1, 2, \ldots, G$:
\begin{align}
  \sigma_{a, g}                         & = 0.25 \, \text{cm}^{-1}, \\
  \sigma_{g \rightarrow g}              & = 1.0 \, \text{cm}^{-1},  \\
  \sigma_{g \rightarrow \parens{g + 1}} & = 1.0 \, \text{cm}^{-1}.
\end{align}

\subsection{Notes from 09/30/2026 meeting}:
\begin{itemize}
  \item \textbf{Need to double-check the following statements/questions}.
  \item With a finite number of gorups, and with downscattering, there are $N$ iterations needed to reach convergence for $N$ groups (?)
  \item Fourier analysis might not be useful here?
  \item Maybe forget about WGS scattering at the moment?
  \item (?) Assume isotropic downscattering
        \begin{itemize}
          \item (?) No scattering source for group $1$, and scattering source
          \item No WGS, only downscatter from group-to-group sequentially (?)
          \item Maybe the solution is going to look like a pure absorber solution?
          \item Maybe make a modification to get rid of residual calculations since we don't need to worry about anisotropic scattering (?)
        \end{itemize}
\end{itemize}

\end{document}
